\Section{Introduction}
\label{sec:intro}

The Move Prover (\MVP)~\cite{DillGPQXZ22} is a formal verification tool for Move smart contracts~\cite{blackshear2020resources,MoveOnAptosBook}, designed for routine use during development with response times comparable to type checkers and linters.
\MVP is deployed at Aptos to verify core protocol logic --- staking, metering, code deployment, and supporting data structures~\cite{ParkZGXGCLC24, AptosFrameworkBook} --- and was originally developed for the Libra/Diem blockchain at Meta.

Recently, Move was extended with \emph{higher-order functions}~\cite{grieskamp2025movehigherorder}: functions are now first-class values that can be passed, returned, stored in struct fields, and persisted in global storage.
Combined with Move's resource model, persisted function values give a principled account of \emph{dynamic dispatch}.
Typical Aptos use cases include callback-based asset flows, extensible asset managers, dispatchable permissions and metering, and generic data-structure operations parameterized by a caller-provided function (e.g.\ fold, map, filter, custom comparators).

Higher-order functions are a known challenge for verification: they break the modular, static-dispatch specification style on which \MVP was built.
When a call is dispatched dynamically, its pre- and post-conditions cannot in general be looked up at the call site; the caller's specification must instead refer \emph{abstractly} to the behavior of its function-value arguments, and the verifier must reconcile this with whichever concrete function reaches the call site.
For a production-oriented prover, this reconciliation must be fast, complete in the relevant cases, and free of the trigger instability (`butterfly' effect) that SMT encodings easily introduce~\cite{BUTTERFLY}.

\Paragraph{Contributions}
This paper describes the extension of \MVP to support higher-order Move, and the underlying extension of the Move specification language:

\begin{Itemize}
    \item We extend Move specifications with \emph{behavioral predicates} $\requiresof{f}{\bar{x}}$, $\abortsof{f}{\bar{x}}$, and $\ensuresof{f}{\bar{x}, \bar{y}}$, which characterize the pre-/abort-/post-behavior of a function value $f$ on arguments $\bar{x}$ and results $\bar{y}$, making function specifications first-class. They implicitly track read and read/write state dependencies.

    \item We introduce \emph{state labels}, which name intermediate memory states along a function's execution and connect them via two-state predicates, as in $\slabeled{S1..S2}{\phi}$, enabling reasoning about sequential evolution of state.

    \item We describe how \MVP implements invocation of function values on top of SMT~\cite{SMTLib2010,z3} by discriminating over the function-value variants flowing into the call site: if $f$ is known, the invocation is replaced by the effect of the concrete function; otherwise (e.g., a function-typed parameter or a closure loaded from storage), the behavioral predicates of $f$ describe its effect.

    \item We extend \MVP's \emph{specification inference} tool \cite{GrieskampEtAl26Inf} to function-valued code.
    The tool computes weakest preconditions over Move bytecode with closures and dynamic dispatch, automatically producing behavioral-predicate-based specifications.
    This both reduces the annotation burden and yields a validation pipeline confirming the encoding is internally consistent.
\end{Itemize}

\noindent While these aspects have similarities to features and solutions in Dafny, \fstar, or \verus, they are novel in the details. For example, Dafny has state labels but does not support higher-order functions with side effects, nor quantification over state labels. See Sec.~\ref{sec:related} for a more thorough discussion of related work.

%Sec.~\ref{sec:move-hof} reviews Move on Aptos, introduces the running AMM example, and illustrates the specification patterns.
%Sec.~\ref{sec:enc} describes the encoding of behavioral predicates, state labels, and dynamic dispatch.
%Sec.~\ref{sec:validation} extends \MVP's specification inference to higher-order code.
%Sec.~\ref{sec:conclusion} discusses related work and concludes.

%%% Local Variables:
%%% mode: latex
%%% TeX-master: "main"
%%% End:
